\section{الخلفية}
\label{sec:background}

تشمل أنظمة أمن المعلومات عمومًا \textbf{التعمية} و\textbf{الخصوصية} و\textbf{الإخفاء}، ويعد الأخير، المعروف باسم \textbf{إخفاء المعلومات}، محور هذه الدراسة. فالتعمية والخصوصية تحميان محتوى الرسالة، لكنهما لا تخفيان وجود الاتصال نفسه، وقد يكون هذا الوجود مثيرًا للريبة. لذلك يعطي إخفاء المعلومات الأولوية لخاصية \textbf{اللامحسوسية}: أي تضمين المعلومات في حوامل عادية، مثل الصور أو النصوص، بحيث تبقى الرسائل المخفية غير ملحوظة.

يمثل النص حاملًا صعبًا على نحو خاص بسبب انخفاض التكرار فيه وصرامة قيوده الدلالية. وتوضح مسألة \textit{السجناء} الكلاسيكية \cite{simmons1984prisoners} الهدف: إذ يجب على طرفين، أليس وبوب، تبادل معلومات مخفية دون تنبيه خصم مراقب.

تقسم طرائق الإخفاء النصي للمعلومات عادة إلى طرق \textbf{قائمة على التنسيق} تستغل خصائص التخطيط أو البنية، وطرق \textbf{لغوية} تعدل الصيغة اللغوية. وضمن الصنف الثاني، تقوم تقنيات مبكرة مثل \textbf{استبدال المرادفات} بتضمين البتات عبر تغيير الاختيارات المعجمية، لكنها تعاني من انخفاض السعة وارتفاع قابلية الكشف.

وبصورة أكثر رسمية، يعمل الإخفاء اللغوي للمعلومات التوليدي من خلال \textbf{دالة تضمين} $f_{emb}$ تنشئ مقابلة تعمية ولغوية لإخفاء الرسائل السرية داخل النص المولد. تأخذ الخوارزمية نموذجًا لغويًا $\mathcal{L}$ ورسالة سرية $m$ ممثلة كسلسلة بتات موزعة بانتظام. وفي كل خطوة توليد، لا تؤخذ العينة بحرية من التوزيع الشرطي لنموذج اللغة $p_{LM}$، بل تقوم دالة التضمين $f_{emb}$ \textbf{بتقييد أو تعديل} اختيار الرمز وفق بتات الرسالة المراد تضمينها. وتعرف هذه العملية باسم \textbf{أخذ العينات الإخفائي}، وتنتج نصًا إخفائيًا $y$. ويعتمد أمن النظام على تقليل التباعد بين توزيع نموذج اللغة الأصلي $p_{LM}$ والتوزيع المعدل $q$ الذي تحدثه $f_{emb}$ \cite{zhang2021provably}.

توفر الطرق اللغوية التقليدية سعة تضمين محدودة وتترك غالبًا آثارًا إحصائية. أما التقدم في التعلم العميق و\textbf{نماذج اللغة الكبيرة} فيتيح الآن طرائق توليدية تحقق جودة نصية أعلى وتضمينًا أكثر أمنًا. ويتطلب تقييم هذه الأنظمة أبعادًا متعددة من اللامحسوسية: \textbf{إدراكية حسية}، أي طبيعية النص للبشر، و\textbf{إحصائية}، أي التشابه التوزيعي مع النص الطبيعي، و\textbf{معرفية}، أي الوفاء الدلالي والسياقي \cite{ding2023context}.

ومع ذلك، يقدم استخدام نماذج اللغة الكبيرة مصادر جديدة لقابلية الكشف. ومن القضايا المهمة \textbf{الهلوسة}، حيث يولد النموذج محتوى غير مدعوم واقعيًا، أو غير متسق دلاليًا، أو ضعيف الارتباط بالسياق السابق. وفي توليد النص العادي، تناقش الهلوسات غالبًا بوصفها أخطاء موثوقية. أما في التوليد الإخفائي، فتغدو كذلك مشكلة أمنية: فحتى إن بدت الجملة فصيحة سطحيًا، فإن ادعاءً غير مدعوم أو انتقالًا مفاجئًا في الموضوع أو تناقضًا مع السياق المحيط قد يجعل النص المُوَرّى مشبوهًا. لذلك تهدد الهلوسة أساسًا \textbf{اللامحسوسية المعرفية}، لأن النص المولد يجب ألا يبدو طبيعيًا فحسب، بل يجب أن يبقى معقولًا سياقيًا ومتماسكًا تواصليًا.

ويتمثل تحد ثان في \textbf{أثر \EN{PSIC}}، أي تعارض اللامحسوسية الإدراكية والإحصائية \cite{yang2020vae}، وهو يجسد توترًا مركزيًا في الإخفاء اللغوي للمعلومات التوليدي. فالنص المحسن للطلاقة البشرية ليس بالضرورة غير قابل للتمييز إحصائيًا عن النص البشري الطبيعي. فعلى سبيل المثال، قد يؤدي اختيار الرموز عالية الاحتمال فقط إلى جمل ذات حيرة منخفضة وقابلية قراءة محلية عالية، لكنه قد يجعل التوزيع المولد أكثر سلاسة وانتظامًا وتركيزًا من الكتابة البشرية الحقيقية. وقد يبدو مثل هذا النص طبيعيًا للقراء، بينما يبقى قابلًا للكشف عبر محللات الإخفاء الإحصائية \cite{yang2020vae}.

لذلك يبين أثر \EN{PSIC} أن \textbf{اللامحسوسية الإدراكية} و\textbf{اللامحسوسية الإحصائية} هدفان مترابطان لكنهما غير متطابقين. فالنص البشري يتضمن غالبًا تنوعًا ولا انتظامًا واختيارات كلمات غير متوقعة محليًا. أما النص المولد شديد الطلاقة فقد يفتقر إلى هذا التباين الطبيعي. ويمكن لزيادة معدل التضمين أحيانًا أن تجبر آلية أخذ العينات على الاختيار من نطاق أوسع من الرموز، بما يقرب النص الناتج من تنوع النص البشري؛ غير أن العملية نفسها قد تقلل قابلية القراءة أو الاتساق الدلالي إذا أصبحت اختيارات الرموز ضعيفة الملاءمة للسياق. ومن ثم فإن السعة الأعلى لا تعني تلقائيًا أمنًا أفضل. يتطلب الإخفاء التوليدي للمعلومات الآمن موازنة معدل التضمين والطلاقة والتشابه التوزيعي والاتساق السياقي.

\subsection{نماذج اللغة الكبيرة بوصفها مولدات نصية}
\label{sec:background-llm-generators}

نماذج اللغة الكبيرة (\EN{LLMs}) أنظمة توليدية ذاتية الانحدار تستند إلى بنية المحول \EN{(Transformer)} \cite{vaswani2017attention}، وتقارب توزيعات عالية الأبعاد على متتاليات اللغة الطبيعية \cite{kaptchuk2021meteor,radford2019language}. فبالنظر إلى سابقة نصية، يصدر النموذج متجه احتمالات على المفردات، ثم تؤخذ عينة الرمز التالي من هذا المتجه وتُلحق بالسابقة، وتتكرر العملية حتى يتحقق معيار توقف. وأثناء التدريب المسبق، تُضبط مليارات المعاملات على مجموعات نصية ضخمة كي يقارب التوزيع التنبؤي للنموذج الأنماط الموجودة في بيانات التدريب \cite{brown2020language}. ونتيجة لذلك، تستطيع نماذج اللغة الكبيرة الحديثة إنتاج نص بطلاقة وتماسك وتحكم أسلوبي عاليين \cite{bubeck2023sparks}. وتلتقط التمثيلات المتعلَّمة انتظامات أسلوبية ودلالية تعمم عبر نطاقات مختلفة، مما يتيح تطبيقات تتطلب محاكاة لغوية دقيقة \cite{reif2022recipe}.

\subsection{نماذج اللغة الكبيرة في الإخفاء اللغوي التوليدي}
\label{sec:background-llm-stego}

بما أن النص الذي تنتجه نماذج اللغة الكبيرة الحديثة يمكن أن يشابه إلى حد كبير النص البشري العادي، فإن هذه النماذج توفر محركات فعالة للإخفاء النصي للمعلومات التوليدي. وبدلًا من تعديل حامل موجود مسبقًا، يعامل عدد متزايد من الأعمال النموذج نفسه بوصفه أداة أخذ عينات إخفائية: إذ توجه الرسالة السرية عملية التوليد، بحيث يُستمد النص المُوَرّى من توزيع تواصلي واقعي منذ البداية. وقد جعل التوفر العام لنماذج عالية الجودة، إلى جانب كفاءة التضمين القائم على أخذ العينات، هذا المسار أكثر عملية من التصاميم السابقة.

تُستخدم نماذج مثل \EN{GPT-2} \cite{radford2019language} و\EN{LLaMA} \cite{touvron2023llama2} و\EN{Baichuan2} \cite{yang2023baichuan2} عادة بوصفها نماذج توليدية أساسية للإخفاء المعلوماتي. وتقرن معظم هذه المخططات نموذج اللغة برمز إخفائي يربط بتات الرسالة باختيارات الرموز: ففي كل خطوة توليد، تحدد البتات المراد إخفاؤها أي رمز مرشح يصدر عن التوزيع التالي للنموذج. غير أن طرائق \textbf{الصندوق الأبيض} التقليدية تستلزم مشاركة نموذج اللغة ومفرداته التدريبية بدقة، مما يحد من الطلاقة والمنطق والتنوع مقارنة بالنصوص التي تولدها واجهات نماذج اللغة الكبيرة الحديثة القائمة على \textbf{الصندوق الأسود}. وتغيّر هذه الطرائق حتمًا توزيعات احتمال أخذ العينات عند تضمين الرسائل، مما يشكل مخاطر أمنية \cite{wu2024generative}. وتتجنب \textbf{طرائق الصندوق الأسود} هذه القيود من خلال الاستفادة من أحدث النماذج المستضافة لتقديم جودة نصية أفضل، وسرعة توليد أعلى، ومتطلبات موارد محلية أدنى. غير أن هذا النهج يستبدل التحكم الدقيق في أخذ العينات الداخلي بهذه المزايا العملية. وعلى النقيض، يتيح \textbf{الوصول إلى الصندوق الأبيض} تحكمًا دقيقًا بالمعاملات وضماناتٍ أمنية أقوى، لكنه يتطلب موارد حاسوبية كبيرة وجهدًا هندسيًا وزمن استجابة أعلى، مما يرفع حواجز النشر. ويُعد هذا التوازن بين نمطي الوصول محوريًا في تقييم متانة الإخفاء اللغوي للمعلومات الحديث وقابليته للتطبيق، وهو ما يحفز مباشرة اختيار تصميم الصندوق الأسود المعتمد في القسم~\ref{sec:methodology}.

وتستكشف طرائق جديدة، مثل \EN{LLM-Stega} \cite{wu2024generative}، الإخفاء النصي للمعلومات التوليدي بنموذج الصندوق الأسود باستخدام واجهات المستخدم لنماذج اللغة الكبيرة، مما يزيل الحاجة إلى مراقبة توزيعات أخذ العينات الداخلية. إذ تُعمَّى بتات الرسالة السرية وتُقابل مع مجموعة كلمات مفتاحية متفق عليها مسبقًا توجه التوليد، وتُستبعد المخرجات المرشحة التي يتعذر استرجاع الرسالة منها عبر أخذ عينات رافض، بما يحافظ على دقة الاسترجاع والجودة الدلالية معًا \cite{wu2024generative}.

ويستفيد إطار آخر، هو \EN{Co-Stega}، من نماذج اللغة الكبيرة لمعالجة تحدي انخفاض السعة في وسائل التواصل الاجتماعي. إذ يوسع الحيز المتاح لإخفاء الرسائل من خلال استرجاع السياق، ويستخدم موجهات مصممة خصيصًا لرفع إنتروبيا الردود المولدة، مما يوسع بدوره سعة التضمين مع الحفاظ على طلاقة المخرجات وملاءمتها الموضوعية \cite{liao2024co}.

أما \textbf{الإخفاء اللغوي بلا لقطات} \EN{(Zero-shot)} فيعتمد بدلًا من ذلك على التعلم في السياق: إذ تُقدَّم عينات نصية حاملة ضمن الموجه، ويُنتج النص المُوَرّى ضمن تبادل سؤال وجواب، بحيث تُقرأ المخرجات بوصفها ردًا طبيعيًا لا نصًا مولدًا قائمًا بذاته \cite{lin2024zero}.

وقد جعل تزايد شيوع النماذج التوليدية العميقة من الممكن تطبيق إخفاء المعلومات القابل للإثبات أمنيًا في سيناريوهات واقعية، إذ تستوفي متطلبات أخذ العينات المثالي والتوزيعات الاحتمالية الصريحة \cite{ding2023discop, kaptchuk2021meteor, qi2024provably}.

وتُقيَّم طرائق إخفاء المعلومات القائمة على نماذج اللغة الكبيرة عادة وفق محورين رئيسين: اللامحسوسية (بأبعادها الإدراكية والإحصائية والمعرفية)، وسعة التضمين (مثل بتات لكل رمز أو بتات لكل كلمة). وقد تشمل تقييمات اللامحسوسية مقاييس آلية (\EN{Perplexity}، \EN{Distinct-n}، \EN{KL/JSD}) إلى جانب أحكام بشرية؛ وتُذكر سعة التضمين عادة كبتات/رمز أو كمعدل تضمين إجمالي. ويستعرض القسم~\ref{sec:related_work} كيفية تطبيق الأدبيات لهذه الطرائق وتقييمها وتقييدها بتفصيل أكبر.
